\section{Q4：Kuramoto 模型与电网——为什么能用、代价是什么（重点）}
\label{sec:q4}

\noindent\textbf{本章回答}：Kuramoto 为什么能用于电网、假设是什么、与 RMS 是什么关系。
逻辑是“先建模型、再证明它从哪来、最后划清边界”：先给标准形式与变量含义，
再从 swing 方程逐步推出它，最后给出可解析结论与适用边界。

\subsection{模型形式与变量含义}

Kuramoto 模型描述 $N$ 个弱耦合、近全同的相位振子\cite{kuramoto1984}：
\begin{equation}
\dot\theta_i=\omega_i+\frac{K}{N}\sum_{j=1}^{N}\sin(\theta_j-\theta_i),
\qquad i=1,\dots,N,
\label{eq:kuramoto}
\end{equation}
\textbf{读法}：每个振子按自身的自然频率 $\omega_i$ 转动，同时被耦合项“拉”向与其他振子一致的相位；
$K$ 越大，拉得越紧。
其中三个量的物理含义是
\begin{center}
\begin{tabular}{@{}ll@{}}
\toprule
符号 & 含义 \\
\midrule
$\theta_i$ & 第 $i$ 个振子的相位（在电网中对应节点电压相角或机组转子角） \\
$\omega_i$ & 自然频率（无耦合时该振子自由振荡的频率；电网中对应节点有功不平衡） \\
$K$ & 耦合强度（电网中对应线路电纳/电抗，$K_{ij}=V_iV_jB_{ij}$） \\
\bottomrule
\end{tabular}
\end{center}
描述宏观同步程度用\textbf{序参量}
\begin{equation}
r e^{\jj\psi}=\frac{1}{N}\sum_{j=1}^{N}e^{\jj\theta_j},\qquad r\in[0,1],
\label{eq:order-param}
\end{equation}
\textbf{读法}：$r=0$ 表示完全无序，$r=1$ 表示完全同步，$\psi$ 为平均相位——
用一个标量概括整个网络的同步程度。
由于式~\eqref{eq:kuramoto} 的耦合项可写成
$\frac{K}{N}\sum_j\sin(\theta_j-\theta_i)
=K\,\Imag\bigl[e^{-\jj\theta_i}\frac{1}{N}\sum_j e^{\jj\theta_j}\bigr]
=K r\sin(\psi-\theta_i)$，
全连接 Kuramoto 可用\textbf{平均场}精确积分（每步只需 $O(N)$）：
\begin{equation}
\dot\theta_i=\omega_i+K\,r(t)\sin\bigl(\psi(t)-\theta_i\bigr).
\label{eq:kuramoto-meanfield}
\end{equation}
\textbf{读法}：耦合被压缩成“与平均场的相互作用”，每步只需 $O(N)$——
这正是 Kuramoto 可解析、可做大规模数值实验的原因。
这正是本报告 \texttt{fig11} 数值实验的计算方式。

\subsection{为什么能用于电网：从 swing 方程到 Kuramoto（假设 A1--A3）}
\begin{derivblock}{推导块 5：从 swing 方程到 Kuramoto}
\textbf{起点}：swing 方程 $+$ 网络功率；\textbf{目标}：三假设下的退化（3 步）；\textbf{结论}：Kuramoto 是 swing 的严格退化；\textbf{支撑}：Q4“为什么能用”。



\paragraph{第一步：转子运动方程。}
同步机的二阶经典模型为\cite{anderson2003,kundur1994}
\begin{equation}
\dot\delta_i=\omega_i-\omega_s,\qquad
M_i\,\dot\omega_i=P_{m,i}-P_{e,i}-D_i(\omega_i-\omega_s),
\label{eq:swing}
\end{equation}
$M_i$ 为惯性常数，$D_i$ 为阻尼，$P_{m,i}$ 为机械功率，$P_{e,i}$ 为电磁功率。
\textbf{读法}：第一式给出相角与转速的关系；第二式是转子运动方程——
不平衡功率（$P_m-P_e$）驱动转速变化，阻尼 $D$ 抑制振荡。

\paragraph{第二步：电磁功率的网络表达。} 由式~\eqref{eq:power-flow}，
节点 $i$ 的有功注入为
\begin{equation}
P_{e,i}=V_i^2G_{ii}+\sum_{j\neq i}V_iV_j\bigl(G_{ij}\cos\delta_{ij}+B_{ij}\sin\delta_{ij}\bigr),
\qquad \delta_{ij}=\delta_i-\delta_j .
\label{eq:Pe-general}
\end{equation}
\textbf{读法}：电磁功率由网络导纳与所有相角差共同决定——
网络约束就是在这里进入动力学方程的。

\paragraph{第三步：三个假设。} 在\textbf{（A1）无损}（$G_{ij}=0$）、
\textbf{（A2）电压幅值恒定}（$V_i\equiv V$）、
\textbf{（A3）均匀惯性/阻尼}（$M_i\equiv M$，$D_i\equiv D$）的条件下
（这一化简在文献\cite{filatrella2008,dorfler2012}中被严格讨论），
式~\eqref{eq:Pe-general} 退化为
\begin{equation}
P_{e,i}=\sum_{j\neq i}K_{ij}\sin\delta_{ij},\qquad K_{ij}=V_iV_jB_{ij}=\frac{V_iV_j}{X_{ij}} .
\label{eq:Pe-lossless}
\end{equation}
\textbf{读法}：在 A1、A2 下只剩 $\sin$ 项；耦合强度 $K_{ij}=V_iV_j/X_{ij}$
完全由电压幅值与线路电抗决定，与运行点无关。
把式~\eqref{eq:Pe-lossless} 代入式~\eqref{eq:swing}：
\begin{equation}
\boxed{\;
M_i\ddot\delta_i=P_{m,i}-\sum_{j\neq i}K_{ij}\sin(\delta_i-\delta_j)-D_i\dot\delta_i\;}
\label{eq:swing-kuramoto}
\end{equation}
式~\eqref{eq:swing-kuramoto} 保留惯性，称为\textbf{二阶 Kuramoto}（含惯性的 swing 型）。
若进一步取\textbf{过阻尼极限} $M_i\to0$，则惯性项消失，化为一阶标准 Kuramoto 式~\eqref{eq:kuramoto}。
因此：
\begin{quote}
\textbf{Kuramoto 模型不是“类比”，而是同步机摇摆方程在（A1）无损、（A2）恒幅值、
（A3）均匀惯性/过阻尼三个假设下的严格退化形式。}
其变量对应为：$\theta_i\leftrightarrow\delta_i$，
$\omega_i\leftrightarrow P_{m,i}/M_i$（或 $P_{m,i}/D_i$），
$K_{ij}\leftrightarrow V_iV_j/X_{ij}$。
\end{quote}

\end{derivblock}

\begin{figure}[htbp]
\centering
\includegraphics[width=0.99\textwidth]{fig10_swing_vs_kuramoto.pdf}
\caption{\textbf{这张图说明}：Kuramoto 与真实的电网模型差多少、在什么条件下完全一样。
\textbf{画的是什么}：(a) $R=0$ 时三台机组的转子角 $\delta_i(t)$（彩色实线为完整 swing 模型，黑色虚线为 Kuramoto 模型，两者重合）；
(b) $R/X=0.5$ 时两者分离（实线为完整模型、虚线为 Kuramoto）；
(c) 稳态相对角偏差（各时刻先减去整体均值再取最大）与 $R/X$ 的双对数关系及斜率 1 参考线。
$R,X$ 为线路电阻与电抗；假设 A1（无损）、A2（恒幅值）见 \S\ref{sec:q4}。
\textbf{怎么读}：(a) 两解完全重合，说明在假设 A1--A2 下 Kuramoto 是严格退化；
(b)(c) 电阻使两解分离，且相对角偏差与 $R/X$ 成正比。
\textbf{直观理解}：Kuramoto 只保留了“纯电抗”的耦合，电阻带来的那一部分被丢掉；
电阻相对电抗越大（$R/X$ 越大），丢掉的越多。
\textbf{结论}：Kuramoto 丢掉的是电导 $G$ 项；输电线路 $R/X\sim0.05$--$0.2$ 可接受，配电网 $R/X\sim1$--$5$ 不适用。}
\label{fig:swing}
\end{figure}

图~\ref{fig:swing} 的数值实验还揭示一个细节：
$R\neq0$ 时系统的整体相角会出现\textbf{中性漂移}（所有 $\delta_i$ 同时平移），
因为式~\eqref{eq:swing-kuramoto} 对整体旋转不变，所有 $\delta_i+c$ 都是平衡点。
因此评价 Kuramoto 的“形状误差”时必须先减去整体均值，
否则会把中性模式误记为模型误差。这一处理已写在脚本 \texttt{fig10\_swing\_kuramoto.py}
的 \texttt{rel\_dev()} 中。

\paragraph{小问题：凭什么说 Kuramoto 能用在电网上？假设有多强？}
不是类比，而是严格退化：摇摆方程式~\eqref{eq:swing}
$+$ 无损（A1）$+$ 恒幅值（A2）$+$ 均匀惯性/过阻尼（A3）
即得式~\eqref{eq:swing-kuramoto}；图~\ref{fig:swing}(a) 给出数值验证（两者轨迹完全重合）。
假设强度可直接量化：放开 A1 后偏差 $\propto R/X$（图~\ref{fig:swing}c），
输电线路 $R/X\sim0.05$--$0.2$ 时偏差在百分之几，配电网 $R/X\sim1$--$5$ 时失效。


\subsection{可解析结论：同步相变与临界耦合}
\begin{derivblock}{推导块 6：自洽方程与临界耦合}
\textbf{起点}：序参量 $+$ 平均场；\textbf{目标}：临界耦合（4 步）；\textbf{结论}：$K_c=2/(\pi g(0))$；\textbf{支撑}：承载力宏观指标。



式~\eqref{eq:kuramoto} 最重要的性质是\textbf{同步相变}\cite{dorfler2013,rohden2012}：
当耦合强度低于临界值 $K_c$ 时
$r\to0$（无序），超过后 $r$ 连续增大。对自然频率分布 $g(\omega)$，
经典结果（Strogatz\cite{strogatz2000,acebron2005}）为
\begin{equation}
K_c=\frac{2}{\pi\,g(0)} .
\label{eq:Kc}
\end{equation}
\textbf{读法}：临界耦合只取决于\textbf{自然频率分布在零点处的密度} $g(0)$——
分布越集中（$g(0)$ 越大），\textbf{门槛 $K_c$ 越低}。
\textbf{方向不要记反}：门槛低意味着“在同样的 $K$ 下已经同步”，
而不是“$K$ 越小越容易同步”。固定 $g(\omega)$ 时，$r$ 随 $K$ 单调增
（$K\le K_c$ 时 $r=0$；$K>K_c$ 后 $r$ 由零连续上升），
对 Lorentzian 分布有闭式 $r=\sqrt{1-K_c/K}$，可作单调性的直接验证。
其推导是：在同步态下，能被锁定的振子满足 $|\omega|\le Kr$ 且相位满足
$\sin\theta=\omega/(Kr)$，故其条件密度为
$\rho(\theta|\omega)=\delta\bigl(\theta-\arcsin\frac{\omega}{Kr}\bigr)$；
漂移振子对序参量的实部贡献为零。代入序参量定义式~\eqref{eq:order-param} 的实部，
得\textbf{自洽方程}
\begin{equation}
r=\int_{-Kr}^{Kr}\sqrt{1-\Bigl(\frac{\omega}{Kr}\Bigr)^2}\,g(\omega)\,\dd\omega .
\label{eq:selfconsistent}
\end{equation}
\textbf{读法}：这是一个\textbf{自洽}方程——右端通过积分上限与分母依赖 $r$，只能迭代求解；
它把“宏观同步程度 $r$”与“微观频率分布 $g(\omega)$”联系了起来。
为求临界耦合，令 $r\to0^+$。做变量替换 $\omega=Kru$，式~\eqref{eq:selfconsistent} 化为
$r=Kr\int_{-1}^{1}\sqrt{1-u^2}\,g(Kru)\,\dd u$；
当 $r\to0$ 时 $g(Kru)\to g(0)$，并利用 $\int_{-1}^{1}\sqrt{1-u^2}\,\dd u=\pi/2$，
两边约去 $r$ 得 $1=K g(0)\pi/2$，即式~\eqref{eq:Kc}。对 Lorentzian 分布 $g(\omega)=\frac{\gamma}{\pi(\omega^2+\gamma^2)}$，
$K_c=2\gamma$。

\end{derivblock}

\begin{figure}[htbp]
\centering
\includegraphics[width=0.99\textwidth]{fig11_kuramoto_transition.pdf}
\caption{\textbf{这张图说明}：耦合要强到什么程度，整个网络才会同步。
\textbf{画的是什么}：(a) 稳态序参量 $r$（$r=0$ 为无序、$r=1$ 为完全同步）与耦合强度 $K$ 的关系（带点线），
红色竖虚线为理论值 $K_c=2\gamma$（$\gamma$ 为自然频率分布 $g(\omega)$ 的 Lorentzian 半宽，本例 $\gamma=1$，故 $K_c=2$）。
(b) 两种耦合强度下 $r$ 随时间的演化 $r(t)$。
(c) 各振子的平均频率 $\langle\dot\theta_i\rangle$ 与自然频率 $\omega_i$ 的关系，黑虚线为未耦合情形 $\dot\theta=\omega$。
振子数 $N=400$。
\textbf{怎么读}：(a) $r$ 在 $K_c$ 处由零转为正，即同步相变；(c) $K>K_c$ 时中部一段振子的频率被锁定（水平段），边缘振子仍漂移。
\textbf{直观理解}：耦合太弱时，每个振子按自己的频率转，谁也带不动谁；
耦合超过一个门槛后，中间的振子被“拉住”、频率趋于一致——这个门槛就是 $K_c$。
门槛高低取决于“大家频率差得有多大”（分布越宽越难同步）。
\textbf{结论}：临界耦合有解析式 $K_c=2/(\pi g(0))$，可作为宏观同步性指标。}
\label{fig:transition}
\end{figure}

电网映射：$K_c$ 越大表示需要越强的网架/越大的耦合才能维持同步；
这为“承载力”提供了一个\textbf{宏观同步性指标}（序参量 $r$）的建模思路。
但必须注意图~\ref{fig:grid} 的结论：拓扑稀疏时同步阈值显著提高，
全连接的式~\eqref{eq:Kc} 不能直接搬到电网。

\begin{figure}[htbp]
\centering
\includegraphics[width=0.99\textwidth]{fig12_kuramoto_grid.pdf}
\caption{\textbf{这张图说明}：在真实网络拓扑上，同步是更容易还是更难。
\textbf{画的是什么}：(a) IEEE 9 节点拓扑，节点颜色按 $\sin\theta_i$ 着色（$\theta_i$ 为节点相角），此例取弱耦合 $K=0.5$。
(b) 三种耦合强度下序参量 $r(t)$ 的演化。
(c) $K=8$ 时各节点相角（减去各自初值）随时间的演化。
\textbf{怎么读}：(a) 颜色分散表示相位分散；(b) 耦合越强 $r$ 越大、越早趋稳；(c) 强耦合时所有曲线合成一根斜线，即相角被锁定。
\textbf{为什么}：耦合项沿网络的边累加，稀疏网络中每个节点只与少数邻居相连，
有效耦合更弱，因此必须用更大的 $K$ 才能达到同样的同步程度；
这也是全连接公式 $K_c=2/(\pi g(0))$ 不能直接外推的原因。
\textbf{直观理解}：耦合是沿着边“传递”的——邻居越少，每个振子受到的拉力越弱，
所以要更大的 $K$ 才能把它们拉成同步。
\textbf{结论}：拓扑稀疏时同步更难，全连接的 $K_c$ 公式不能直接搬到真实网络。}
\label{fig:grid}
\end{figure}

\paragraph{小问题：$K_c=2\gamma$ 对电网意味着什么？}
$\gamma$ 是节点有功不平衡（自然频率）分布的离散程度，$K_c$ 是维持同步所需的
最小耦合强度。映射到电网：
$K_c$ 越大，说明需要越强的网架（$B_{ij}=1/X_{ij}$ 越大）或越小的不平衡才能同步。
于是“接入更多 DER”在 Kuramoto 语言里等价于“$\gamma$ 增大而 $K_c$ 不变”，
结果是同步裕度下降。


\subsection{与 RMS 的关系：降阶，而非等价}

综合 \S\ref{sec:singular} 与本节推导，四类模型之间是一个\textbf{简化关系}——注意其中
PF 与 Kuramoto 是两条并列的分支，而不是前后两级：

\begin{equation}
\begin{array}{rcl}
\text{EMT} & \xrightarrow[\text{平均开关}]{\text{忽略网络电磁暂态}} & \text{RMS}\\[5pt]
\text{RMS} & \xrightarrow[\text{令所有导数}=0]{\ \ \text{静态化}\ \ } & \text{PF / QSTS}\\[5pt]
\text{RMS} & \xrightarrow[\text{恒幅值、无损纯感}]{\text{相角化（消去网络）}} & \text{Kuramoto}
\end{array}
\label{eq:chain}
\end{equation}
\textbf{读法}：前两行是一条链（EMT $\to$ RMS）；后两行是 RMS 的\textbf{两个并列简化}——
PF 丢掉时间维度（静态化），Kuramoto 保留时间但消去网络与电压（相角化）。
因此 Kuramoto \textbf{不是}“潮流的进一步简化”：它比 PF 保留了转子摇摆动态，
只是把状态从“电压 $+$ 相角”压缩到“仅相角”。
表~\ref{tab:kuramoto-rms} 给出 Kuramoto 与 RMS 的逐项对照。

\begin{table}[htbp]
\centering
\caption{Kuramoto 与 RMS 的关系}
\label{tab:kuramoto-rms}
\small
\begin{tabularx}{\textwidth}{@{}lXX@{}}
\toprule
 & \textbf{RMS} & \textbf{Kuramoto} \\
\midrule
状态变量 & $\delta,\omega,E_q',E_d'$, 控制器状态 & 仅 $\delta_i$（一阶）或 $\delta_i,\omega_i$（二阶） \\
电压幅值 & 变量（$V_i$ 动态由网络/励磁决定） & 常数 $V_i\equiv V$（假设 A2） \\
网络 & 代数约束式~\eqref{eq:power-flow}，含 $G_{ij},B_{ij}$ & 直接消去为 $\sin$ 耦合，仅 $B_{ij}$（假设 A1） \\
无功/损耗 & 保留 & 全部丢弃 \\
可解析性 & 一般需数值积分 & 有相变等解析结论（式~\eqref{eq:Kc}） \\
适用场景 & 通用（输电/配电、弱网） & $X/R\gg1$、强电压支撑、宏观同步问题 \\
\bottomrule
\end{tabularx}
\end{table}

\paragraph{结论。}
Kuramoto 是 RMS 在三个强假设下的\textbf{降阶模型}，不是等价模型：
\begin{enumerate}[leftmargin=1.6em]
\item 在高压输电网（$X/R\sim5\text{--}20$、电压支撑强）的\textbf{同步稳定性}问题上，
它保留了主要机制，可作为解析工具（相变、临界耦合）；
\item 在配电网（$R/X\sim1\text{--}5$）和以电压问题为主导的场景中，
假设 A1--A2 均不成立（图~\ref{fig:swing}c 已量化 $R/X$ 的影响），
Kuramoto 只能做定性/宏观指标，不能替代潮流或 RMS；
\item 对本项目，Kuramoto 的价值在于为“城市电网承载力的宏观同步性判据”
提供一个\textbf{可解析、可解释}的骨架，但必须同时报告其假设边界。
此外，变流器下垂控制（droop control）与 Kuramoto 已有严格对应\cite{simpsonporco2013}，
可作为 DER 占比进一步提高后的扩展方向。
\end{enumerate}

\paragraph{小问题：Kuramoto 与 RMS 是等价、近似还是降阶？配电网能用吗？}
是\textbf{降阶}（reduction）而非等价：丢掉了电压幅值动态、无功与网络损耗，
逐项对照见表~\ref{tab:kuramoto-rms}。
配电网（$R/X\sim1$--$5$）中假设 A1 直接失效，且配电网的核心问题通常是
电压越限与三相不平衡而非功角同步，因此 Kuramoto 只能用作定性/宏观指标，
不能替代潮流或 RMS；这一点由图~\ref{fig:swing}(c) 的 $R/X$ 扫描直接量化。

\begin{figure}[htbp]
\centering
\includegraphics[width=0.99\textwidth]{fig13_model_chain.pdf}
\caption{\textbf{这张图说明}：这几类模型到底是什么关系（是本报告的总览图）。
\textbf{看图}：四列方框依次为完整 EMT、RMS、潮流/功率平衡、Kuramoto；
框间箭头标注各次降阶丢掉的动态，下方灰框标注代价或损失。
\textbf{为什么}：PF 是“令所有导数为零”（丢掉时间维度）；Kuramoto 是“保留时间、但用恒幅值 $+$ 无损
把网络代数约束解析消去”（丢掉电压与网络）。两者丢掉的量不同，因此是并列而非前后两级。
\textbf{直观理解}：PF 是把“时间”整个丢掉（只看稳态）；Kuramoto 是保留时间、但把“电压与网络”丢掉。
两者丢掉的不是同一样东西，所以是并列的两条路，而不是前后两级。
\textbf{结论}：EMT、RMS、PF 是一条链，而 Kuramoto 是 RMS 的\textbf{并列简化}（不是链上下一级）；
Q3 回答前两级的连接，Q4 回答 Kuramoto 从哪来、代价是什么。}
\label{fig:chain}
\end{figure}
